
%COVID cancelled standardized testing in the UK and lead to record number of  top grades getting issues for two years straight, and overstreching universities capacities.
How would teachers and schools grade students differently compared to a standardized grading system where teachers have no discretion? The COVID-19 pandemic canceled standardized exams in the UK for two consecutive years, and the final grades were replaced by teachers’ predicted grades. These teacher-assigned grades determined students’ placement outcomes, as the UK higher education system bases admission decisions on applicants’ standardized exam results. The policy led to a 12-percentage-point increase in the number of top grades awarded in 2020, and an even larger increase in 2021, despite the disruptions in learning for the 2021 cohort. Top universities over-accepted students as a majority of their applicants passed the grade requirements, while lower-ranked universities suffered from students shortages. This episode underscores how a change in the grading system reshaped the grade distribution, with far-reaching consequences for the higher education market.

% Why this is economics? Because its localization of information.
I use the complete transfer of grading discretion from a centralized examiner to local teachers to study the difference across schools in their grading leniency, as well as their grading patterns across students with varying backgrounds\footnote{I use the term grading leniency as the institutional difference in the students scores in the standardized exam and the teachers assigned grade, while I use grade bias as the difference in the grades assigned across students with different demographic attributes while controlling for their academic ability.}. Several studies document differences across schools in their grading patterns when teachers are allowed to use their discretion in grading(\cite{10.1162/00335530360698441}, \cite{figlio2006}, \cite{dee2019}), while others show how discretionary grading may allow teachers or schools to reflect their innate biases onto the students' grades. Notable channels includes gender (\cite{lavy2009}, \cite{10.1093/qje/qjz008}), ethnicity ( \cite{burgess2013}, \cite{Alesina2024}), and economic affluence (\cite{hanna2012}). However, less in known about how the differences across schools in their institutional details explains the degree of grading discretion that teachers exhibit when given a chance to incorporate their knowledge of the students into their grades \footnote{\cite{figlio2006} examines different responses from schools to external incentives to increase student performance, finding how private schools mis-reports their students disability status to hide the bad performance of lower performing students. However they do not focus on grades of the students as the schools in their setting used standardized tests to measure a schools performance.}. The lack of evidence accrues from the difficulties for researchers to establish a comparable measure of grading biases across schools, as schools often vary in their methods for measuring their pupil's academic trajectories, as well as the prohibitive costs in running experiments that involves schools at a national scale.

% How schools inflate differently. Static patterns as well as their dynamics.
Using administrative data on final standardized exam  (e.g. \emph{A-levels})  outcomes for nearly 4,000 schools in the UK, I estimate the difference in the grading policies between schools by comparing the extent of disruption in the trajectories of the final grades at each school. A primary identification threat is that cross-school divergences in grades may reflect pre-existing trends rather than the change in grading method. I rule this out by leveraging the historical centralized grading method employed by the A-level examiners (\emph{Comparable Standards approach}); the exam administrators of the A-levels drew the grade boundaries for the A-level so that the overall grade distribution remains constant overtime after conditioning on the academic performance of the test-taking student population. I confirm the stationary property of the grade distribution through multiple stationarity tests, and show that the conditional expected grade of a student follow the same levels overtime. The stationarity property of the expected grade holds even after sub-sampling the sample by school types, students parental income class, and other co-variates, which establishes how the A-levels are designed to balance the difficulty of the A-levels across years by using the test-takers previous academic achievements in a standardized test (e.g.\emph{GCSE}). The annual alignment of student population and the grade distribution allows me to capture the effect of the transfer of grading discretion to local teacher by measuring the deviance between a students' assigned grade to the expected grade that the students would have received if they took the standardized exam. 

My approach for detecting grading leniency across schools resembles the approaches used in \cite{neal2008}; \cite{dee2019}, in which researchers leverage variation in the grading methods used across periods, and estimate the difference in grades between students with academically or demographically similar attributes in different periods\footnote{My approach follows a similar set up as the pandemic completely transferred the grading authority from central examiners to local teachers for the cohorts taking the A-levels in during the pandemic years. }\footnote{Other papers detect anomalies by comparing students near the cutoff of a passing grade (\cite{neal2008}, \cite{diamond}). Grading differences across demographic groups are estimated by comparing the size of the grading discontinuities across the passing threshold by demographic groups. Such approaches do not compare across exams, and derives grading anomalies using variation across the same test.}. I add to the literature by demonstrating how conventional grading methods in standardized tests can be leveraged to estimate grading biases. The standardized tests in the UK mechanically sets the grade boundaries by using variables (available to the researchers), therefore allowing me to rule out any influence of confounders that could potentially bias the assigned grades. My method additionally avoids other potential sources of bias, such as selection into treatment as my identifying variation accrues over the time dimension and ruling out potential sorting  of students into treatment groups. The mechanical re-balancing of students from diverse background to identical scales of achievements in the final grades can be readily confirmed in traditional stationarity tests (\cite{https://doi.org/10.3982/ECTA19367}).

\vspace{0.5em}

I find substantial heterogeneity in how schools loosened their grading policies, with higher degrees of inflation concentrated in privately funded schools and schools with historically lower average performance in standardized A-level exams. I find a nearly 13 percentage point gap in the assignment of top grades between schools in the bottom 10th percentile of quality and those in the top 10th percentile. Privately funded schools (e.g. Independent schools) exhibited greater grade inflation across all levels of school quality, consistently surpassing the inflation observed in publicly funded schools of similar quality. Examining the channels of grade inflation, I find that differences in overall inflation across schools were largely explained by inflation in non-quantitative subject areas. This reflects the innate difference between subjects in how much grading discretion are passed to the graders. I find a 20 percentage point gap between schools in the bottom 10th percentile and those in the top 10th percentile in the awarding of top grades in non-quantitative subjects, with little to no variation across schools in the extent of grade inflation in quantitative subjects. Most schools maintained consistent levels of grade inflation across the two consecutive years, but I observe a downward revision of grading standards in 2021 among independent schools and schools that had stricter grading policies in the initial year.

Turning to grading bias, I find a persistent ethnicity and gender gap in the degrees of grade inflation, as well as its' monotonic relationship with students economic and academic affluence. Conditional on prior performance, female and white students received top grades by nearly 1.4 and 0.5 percentage points respectively. In particular, the gender bias was stronger at historically high performing schools and in quantitative subjects. In 2021, the gender gap increased at schools at extensive margins of the schools quality distribution, and the gender gap across lower performing schools remained negligible. Similarly, I find grading bias against ethnic minorities, as white students received top grades with higher chance than any other ethnic groups. On affluence measures, I find higher degrees of inflation at students with parents in high-income occupations, and academically well-performing students. Notably, the positive relationship between inflation and both affluence measures were more pronounced in lower performing schools. Such results highlights how demographic and socio-economic attributes of students shaped the shift in the grade distribution, but more importantly how the pre-existing sorting and subject choice patterns between students and schools constrained the extent and patterns of grade inflation that schools displayed over their pupils. 

My results contributes to the large empirical literature on grading biases by demonstrating the importance of institutional structure of schools on understanding the grading patterns of their teachers. Previous literature on gender biases find mixed results with respect on the magnitude of the bias, as \cite{hanna2012} and \cite{diamond} finds no evidence and \cite{lavy2018} and \cite{10.1093/qje/qjz008} finds evidence and remedies. My results provides an empirical explanations on how institutional details of schools may account for the discrepancy in the findings. I find a concentration of gender gaps in low performing schools, particularly in their non-quantitative subject groups. I do not find any evidence of grading bias at high performing schools in non-quantitative subjects, which matches papers reporting no-evidence. Regarding ethnicity, \cite{burgess2013} and \cite{Alesina2024}
finds a grading bias against ethnic minorities, which I complement by showing ethnicity based gaps existed across a wide range of schools. Lastly, my results show how grade inflation were limited in schools with historically higher average academic performance. The results demonstrates how grading biases cannot be observed in higher performing schools as the traditional grading measures truncate students ability upwards, and leaves little space for teachers to additionally signal their worth by exercising their grading discretion. 

My results also contributes on the effect of admission policies on university diversity. Test optional policies remove the requirements for students to submit a standardized test. Advocates cite how test-optional policies allows universities to draw applicants from a broader range of students (\cite{10.1257/aer.20231407}), but opponents cites the access to good grades do not necessarily translate into enrollment, especially for students from disadvantaged backgrounds (\cite{NBERw18586})\footnote{Recent empirical work point to the inefficiency of such policies, citing enrollment reduction(\cite{NBERw34260}) and potential crowding out of good grades from  (\cite{NBERw33389}) of disadvantaged students. }. My paper shows how the removal of standardized grade in the UK evened the grade outcome across students with contrasting backgrounds, but fwith limited impact on application success of disadvantaged student into selective universities. The grade inflation mainly assisted conventional incumbent pool of applicants, especially those on the brink of acceptance, in securing their placements. The inflation therefore led to exacerbation of parental income inequality across universities, as their students composition more reflected the composition of their applicant pool. 

The paper is organized as follows: Section 2 provides the institutional background of the higher education system in the UK as well as the details of the policy set in place during the pandemic. Section 3 describes the data source I use in the analysis. Section 4 presents the empirical model and discuss the validity of the estimates along with the institutional details. Section 5 presents the results on heterogeneity of grading patterns across schools, as well as student level patterns, the dynamics of grading patterns, as well as their influence in placing students to their top choices. Section 6 discuss the magnitudes of the results with the literature on grading biases. Section 7 concludes with policy discussion on how my finding can inform policy makers with thinking about test-optional policies. 